\documentclass[12pt]{article}
\usepackage[utf8]{inputenc}
\usepackage[T1]{fontenc}
\usepackage[margin=1in]{geometry}
\usepackage{setspace}
\usepackage{xcolor}
\usepackage[colorlinks=true, urlcolor=blue!60!black, linkcolor=blue!60!black]{hyperref}
\setlength{\parindent}{0em}
\setlength{\parskip}{0.6em}
\renewcommand{\baselinestretch}{1.0}
\begin{document}
\begin{center}
{\LARGE \textit{Jackson v. Metropolitan Edison Co.}, 419 U.S. 345 (1974)\par}
\vspace{0.5em}
{\large Case Brief --- Two-Page Edition\par}
\end{center}

\section{Memory Jogger}
A regulated private utility's shutoff without notice or hearing was not state action; the Court required a close nexus to the specific act and an exclusively public function (\href{https://supreme.justia.com/cases/federal/us/419/345/}{419 U.S.\ 345}).

\section{Detailed Case Facts}
Metropolitan Edison, a privately owned, PUC-regulated utility, disconnected Catherine Jackson's electricity in York, Pennsylvania, after a dispute over unpaid bills and alleged meter tampering, acting under a termination rule in its PUC-filed tariff and without further notice.

\section{Procedural History}
Jackson sued under \S\ 1983 for denial of due process. The district court dismissed for lack of state action; the Third Circuit affirmed.

\section{Judicial Votes}
Affirmed, 6--3. Rehnquist wrote for Burger, Stewart, White, Blackmun, and Powell; Douglas, Brennan, and Marshall each dissented.

\section{Holding}
State action by a regulated private entity requires a ``sufficiently close nexus'' between the State and the specific challenged act. The public-function doctrine covers only powers ``traditionally exclusively reserved to the State,'' not utility service.

\section{Analysis of Opinions}
Rehnquist held that heavy regulation, monopoly status, and the essential nature of the service were insufficient, and that the PUC's passive acceptance of the tariff did not put the State's weight behind the shutoff, unlike the commission's endorsement in \textit{Pollak}. The State was not ``a partner or even a joint venturer.'' Douglas would weigh the factors in the aggregate; Brennan found no controversy; Marshall stressed the state-sanctioned monopoly and warned that the rule would seemingly shield a racial refusal of service.

\section{Examples: Future Applications}
\textbf{Same-Side:} when a nursing home funded by Medicaid transfers a resident on its doctors' judgment, the decision is private (\textit{Blum}). \textbf{Opposite-Side:} a utility condemning an easement under delegated eminent domain is a state actor. \textbf{Fence-Sitter:} a state-funded charter school's expulsion of a student.

\section{Critique}
Critics fault the seriatim analysis and the treatment of commission silence as neutral. \textit{Memphis Light v.\ Craft} required due process for a municipal utility's shutoff, so protection turns solely on ownership.

\section{Key Quotations}
\begin{quote}
``But the inquiry must be whether there is a sufficiently close nexus between the State and the challenged action of the regulated entity so that the action of the latter may be fairly treated as that of the State itself.''
\end{quote}
\noindent --- Justice Rehnquist, for the Court, 419 U.S.\ at 351.

\end{document}
